\chapter{Glosario}
\label{ap:glosario}
\begin{description}[style=nextline,leftmargin=0.5cm]
\item[AUC (ROC-AUC)] Probabilidad de que el modelo califique a un caso positivo (cae)
por encima de uno negativo. 0.5 es azar; 1 es perfecto.
\item[Brier] Error cuadrático medio de una probabilidad predicha. Menor es mejor.
\item[Calibración] Que las probabilidades predichas coincidan con las frecuencias
observadas.
\item[Cascada] Secuencia de impagos donde cada uno provoca otros.
\item[Choque] Evento externo que golpea a una entidad (retiro de liquidez, caída
operativa, quiebra).
\item[Coherencia local $r_i$] Qué tan en fase le llega el ingreso a una entidad
respecto a su propio ciclo ($1$: en fase; $-1$: en contrafase).
\item[Colchón (umbral)] Déficit máximo que una entidad soporta antes del impago.
\item[Control nulo] Mismo experimento con el mecanismo apagado; si el efecto aparece
ahí, es un artefacto.
\item[Déficit] Liquidez que falta para cubrir los pagos del momento.
\item[Dispersa (matriz)] Matriz que solo guarda los valores distintos de cero.
\item[Fase $\theta$] Posición dentro del ciclo de pagos, como la aguja de un reloj.
\item[Fuga de información] Cuando el modelo usa, sin querer, datos que no tendría en
la realidad; infla los resultados.
\item[GroupKFold] Validación cruzada que nunca pone el mismo grupo (red o mundo) en
entrenamiento y prueba a la vez.
\item[Impago] No pagar lo que se debía.
\item[Índice de amplificación $R$] $\beta\rho(W)/\alpha$: menor que 1, la red
amortigua; mayor que 1, amplifica.
\item[Informalidad] Trabajo sin empleador formal; aquí, personas sin capa de nómina.
\item[Kuramoto] Modelo de osciladores acoplados que explica la sincronización.
\item[Laplaciano] Matriz que describe la difusión en una red; conserva el total.
\item[Liquidez] Dinero disponible en el momento en que se necesita.
\item[Parquet] Formato de archivo para tablas.
\item[PD] Probabilidad de impago (\emph{probability of default}).
\item[Radio espectral $\rho(W)$] El mayor valor propio de $W$.
\item[RK4] Método de Runge–Kutta de orden 4 para integrar ecuaciones diferenciales.
\item[Riesgo crediticio] Riesgo de que un deudor no pague su crédito.
\item[Riesgo sistémico] Riesgo de que fallas locales se propaguen a todo el sistema.
\item[Score crediticio] Modelo que estima la PD a partir de variables del cliente.
\item[Semilla] Número que fija el azar: misma semilla, mismo resultado.
\item[SHAP] Método que reparte una predicción entre las variables que la causan.
\item[Solvencia] Tener más activos que deudas en total.
\item[Verdad del terreno] Valores reales (aquí, simulados) que sirven para evaluar
pero que ningún modelo puede ver.
\item[XGBoost] Modelo de muchos árboles de decisión que se corrigen unos a otros.
\item[YAML] Formato de texto para configuraciones.
\end{description}

\chapter{Referencias}
\begin{itemize}[leftmargin=*]
\item Acebrón, J. A., et al. (2005). The Kuramoto model. \emph{Rev. Mod. Phys.} 77, 137.
\item Acemoglu, D., Ozdaglar, A., \& Tahbaz-Salehi, A. (2015). Systemic risk and
stability in financial networks. \emph{American Economic Review} 105(2).
\item Allen, F., \& Gale, D. (2000). Financial contagion. \emph{J. Political Economy} 108(1).
\item Arenas, A., et al. (2008). Synchronization in complex networks. \emph{Physics Reports} 469.
\item Barabási, A.-L. (2016). \emph{Network Science}. Cambridge University Press.
\item Barabási, A.-L., \& Albert, R. (1999). Emergence of scaling in random networks.
\emph{Science} 286.
\item Battiston, S., et al. (2012). DebtRank: Too central to fail? \emph{Scientific Reports} 2.
\item Battiston, S., et al. (2007). Credit chains and bankruptcy propagation in
production networks. \emph{J. Econ. Dyn. Control} 31(6).
\item Boccaletti, S., et al. (2006). Complex networks: structure and dynamics.
\emph{Physics Reports} 424.
\item Boccaletti, S., et al. (2014). The structure and dynamics of multilayer
networks. \emph{Physics Reports} 544.
\item Chung, F. (1997). \emph{Spectral Graph Theory}. AMS.
\item Eisenberg, L., \& Noe, T. H. (2001). Systemic risk in financial systems.
\emph{Management Science} 47(2).
\item Gai, P., \& Kapadia, S. (2010). Contagion in financial networks. \emph{Proc. R. Soc. A} 466.
\item Goodfellow, I., Bengio, Y., \& Courville, A. (2016). \emph{Deep Learning}. MIT Press.
\item Haldane, A. G., \& May, R. M. (2011). Systemic risk in banking ecosystems.
\emph{Nature} 469.
\item Karniadakis, G. E., et al. (2021). Physics-informed machine learning.
\emph{Nature Reviews Physics} 3.
\item Kipf, T. N., \& Welling, M. (2017). Semi-supervised classification with graph
convolutional networks. \emph{ICLR}.
\item Letizia, E., \& Lillo, F. (2019). Corporate payments networks and credit risk
rating. \emph{EPJ Data Science} 8, 21.
\item Merton, R. C. (1974). On the pricing of corporate debt. \emph{J. Finance} 29(2).
\item Newman, M. E. J. (2003). The structure and function of complex networks.
\emph{SIAM Review} 45.
\item Newman, M. E. J. (2010). \emph{Networks: An Introduction}. Oxford University Press.
\item Óskarsdóttir, M., \& Bravo, C. (2021). Multilayer network analysis for improved
credit risk prediction. \emph{Omega} 105.
\item Pikovsky, A., Rosenblum, M., \& Kurths, J. (2001). \emph{Synchronization}.
Cambridge University Press.
\item Raissi, M., Perdikaris, P., \& Karniadakis, G. E. (2019). Physics-informed neural
networks. \emph{J. Comput. Phys.} 378.
\item Rodrigues, F. A., et al. (2016). The Kuramoto model in complex networks.
\emph{Physics Reports} 610.
\item Shi, J., \& Malik, J. (2000). Normalized cuts and image segmentation.
\emph{IEEE TPAMI} 22(8).
\item Tobback, E., et al. (2017). Bankruptcy prediction for SMEs using relational
data. \emph{Decision Support Systems} 102.
\item Watts, D. J. (2002). A simple model of global cascades on random networks.
\emph{PNAS} 99(9).
\item Banco Central do Brasil, Working Papers 551 y 584.
\end{itemize}
\begin{nota}
Las referencias de Tobback et al., Óskarsdóttir y Bravo y Battiston et al.\ (2007) se
citaron de memoria: verifique volumen y páginas antes de usarlas en un documento
formal.
\end{nota}
